\chapter{Server-Side Integration}

\section{Conversation State Management}
\begin{itemize}
    \item Describe the conversation service and its bounded conversation history.
    \item Explain storage of original text, model text, detected language, output language, and translation status.
    \item Explain why conversation state is separate from scene memory.
\end{itemize}

\section{Prompt Composition from Scene Memory}
\begin{itemize}
    \item Describe how the chat service serializes the current world state into textual context.
    \item Explain object ordering by salience, with emphasis on socially relevant people.
    \item Describe the use of latest caption and recent captions as additional grounding context.
    \item Explain how the current implementation realizes grounded retrieval by rendering structured memory into prompts.
\end{itemize}

\section{General Scene-Aware Chat}
\begin{itemize}
    \item Describe the normal chat endpoint for open questions about the current scene.
    \item Explain language normalization, history inclusion, prompt rendering, and provider-agnostic text generation.
    \item Discuss expected question types and answer scope.
\end{itemize}

\section{Object-Focused Dialogue}
\begin{itemize}
    \item Present the object-chat mode as one of the strongest concrete contributions of the current implementation.
    \item Explain object matching by label, ranking by salience, and extraction of per-object attributes and relations.
    \item Describe the crop-caption fallback used when an object has weak structured memory but a tracked crop is available.
    \item Explain how matched object IDs are returned as traceability metadata.
\end{itemize}

\section{Direct Vision Chat}
\begin{itemize}
    \item Describe the route that bypasses scene-memory grounding and sends an image directly to a multimodal model together with a question.
    \item Explain how this mode complements the memory-grounded chat mode and when it is useful.
    \item Discuss its trade-off: richer direct visual reasoning but weaker persistent identity tracking.
\end{itemize}

\section{Multilingual Handling}
\begin{itemize}
    \item Explain the translation layer and the use of canonical model-facing language when needed.
    \item Describe the current English/Czech output modes and how translation is applied to both user input and model response.
    \item Discuss translation-related failure modes and semantic drift.
\end{itemize}

\section{Failure Handling and Practical Limitations}
\begin{itemize}
    \item Describe response fallback behavior under timeouts, malformed outputs, unavailable services, or missing visual context.
    \item Discuss remaining limits of prompt-only grounding, short conversation memory, and scene-memory incompleteness.
    \item Explicitly identify which limitations belong to the current implementation and which are reserved for future work.
\end{itemize}


\section{Role of the \texttt{research} Package}
\begin{itemize}
    \item Explain that the repository contains both production code and research tooling.
    \item Clearly separate the newer structured experiment pipeline from older exploratory scripts and notebooks.
    \item State which parts directly support the thesis evaluation.
\end{itemize}

\section{Structured Experiment Pipeline}
\begin{itemize}
    \item Present the newer \texttt{research.experiments} pipeline as the primary evaluation framework.
    \item Describe its phases in execution order:
    \begin{itemize}
        \item image description generation;
        \item vocabulary mining;
        \item draft scene graph generation;
        \item context-rot experiments;
        \item scene graph evaluation.
    \end{itemize}
    \item Explain the purpose of run metadata, stage metrics, and artifact saving.
\end{itemize}

\section{Reuse of Production Components in Experiments}
\begin{itemize}
    \item Describe the research adapters that import and reuse production detection, caption, LLM, and VLM code.
    \item Explain why this is important scientifically: the offline experiments should evaluate the same or very similar components as the deployed system.
\end{itemize}

\section{Description Generation and Vocabulary Mining}
\begin{itemize}
    \item Describe the description phase: detect objects, then generate detailed captions focused on those detections.
    \item Describe the vocabulary phase: extract candidate predicates and attributes per image, then consolidate them into a final vocabulary.
    \item Explain why vocabulary design is a key variable for scene graph quality and dialogue usefulness.
\end{itemize}

\section{Draft Scene Graph Generation}
\begin{itemize}
    \item Describe the construction of SoM-marked images for offline scene-graph prompting.
    \item Explain prompt rendering from detected objects, vocabulary, and captions.
    \item Describe the use of structured output schemas for draft graph generation.
\end{itemize}

\section{Context-Rot and Sensitivity Experiments}
\begin{itemize}
    \item Explain the idea behind context-rot: reducing available vocabulary and measuring degradation in graph output.
    \item Describe how vocabulary slices are built and evaluated.
    \item Explain how prompt sensitivity is later summarized across runs.
\end{itemize}

\section{Scene Graph Evaluation Metrics}
\begin{itemize}
    \item Define the core metrics implemented by the code:
    \begin{itemize}
        \item strict triplet precision, recall, and F1;
        \item binary relation triplet metrics;
        \item unary attribute metrics;
        \item relation-pair metrics;
        \item predicate-only metrics;
        \item overgeneration and undergeneration rates;
        \item optional normalized graph edit distance.
    \end{itemize}
    \item Explain micro-averaging and per-predicate summaries.
    \item Describe image-potency statistics and why they help characterize scene complexity.
\end{itemize}

\section{System-Level Evaluation Protocol}
\begin{itemize}
    \item Define what should be evaluated on the deployed system itself:
    \begin{itemize}
        \item latency of captioning, detection, memory update, scene graph generation, and end-to-end chat;
        \item stability of tracked identities across frames and scans;
        \item usefulness of Pepper social fusion for person-centered dialogue;
        \item qualitative quality of grounded robot responses in interactive scenarios.
    \end{itemize}
    \item Explain how pipeline stage timings exposed by the server can be used for runtime profiling.
\end{itemize}

\section{Legacy Exploratory and Training Pipelines}
\begin{itemize}
    \item Describe the older exploratory workflows, but keep them clearly secondary to the main evaluation story.
    \item Mention:
    \begin{itemize}
        \item automated SoM scene-graph dataset generation with Grounding DINO and an LLM labeler;
        \item knowledge distillation for detector training data;
        \item detector training for RT-DETR or YOLO;
        \item LoRA-based VLM fine-tuning with Unsloth.
    \end{itemize}
    \item Explain whether these are reported as completed experiments, supporting tooling, or future-ready extensions.
\end{itemize}
